% 第5章图：三个小图及其宽度最小的树分解。
% 树（边12,23,24）：{1,2}-{2,3}-{2,4}，树宽1；四环1-2-3-4-1：{1,2,3}-{1,3,4}，树宽2；完全图K4：单个袋{1,2,3,4}，树宽3。
\begin{tikzpicture}[x=1mm,y=1mm,>=Stealth,
  v/.style={circle,draw=FigureStroke,fill=FigureBlueFill,line width=.7pt,minimum size=5.4mm,inner sep=0pt},
  e/.style={FigureStroke,line width=.7pt},
  bag/.style={rectangle,rounded corners=1.8pt,draw=FigureStroke,fill=white,line width=.7pt,minimum width=11mm,minimum height=6mm,inner sep=1.5pt},
  sc/.style={font=\figurefont\footnotesize,text=FigureInk},
  itf/.style={font=\figurefont\footnotesize,text=FigureGreen},
  note/.style={font=\FigureNoteFont\fontsize{7.5}{9}\selectfont,text=FigureMuted}]
% (a) tree
\begin{scope}
 \node[v] (a1) at (0,30) {$1$}; \node[v] (a2) at (10,30) {$2$};
 \node[v] (a3) at (20,36) {$3$}; \node[v] (a4) at (20,24) {$4$};
 \draw[e] (a1)--(a2) (a2)--(a3) (a2)--(a4);
 \node[bag] (A1) at (10,12) {$\{1,2\}$};
 \node[bag] (A2) at (10,0) {$\{2,3\}$};
 \node[bag] (A3) at (10,-12) {$\{2,4\}$};
 \draw[FigureStroke,line width=.8pt] (A1)--(A2) node[midway,right=1mm,itf]{$\{2\}$};
 \draw[FigureStroke,line width=.8pt] (A2)--(A3) node[midway,right=1mm,itf]{$\{2\}$};
 \node[sc] at (10,-22) {(a) 树：树宽$1$};
\end{scope}
% (b) four-cycle
\begin{scope}[xshift=38mm]
 \node[v] (b1) at (0,36) {$1$}; \node[v] (b2) at (12,36) {$2$};
 \node[v] (b3) at (12,24) {$3$}; \node[v] (b4) at (0,24) {$4$};
 \draw[e] (b1)--(b2)--(b3)--(b4)--(b1);
 \node[bag] (B1) at (6,9) {$\{1,2,3\}$};
 \node[bag] (B2) at (6,-6) {$\{1,3,4\}$};
 \draw[FigureStroke,line width=.8pt] (B1)--(B2) node[midway,right=1mm,itf]{$\{1,3\}$};
 \node[sc] at (6,-22) {(b) 四环：树宽$2$};
\end{scope}
% (c) complete graph K4
\begin{scope}[xshift=74mm]
 \node[v] (c1) at (0,36) {$1$}; \node[v] (c2) at (12,36) {$2$};
 \node[v] (c3) at (12,24) {$3$}; \node[v] (c4) at (0,24) {$4$};
 \draw[e] (c1)--(c2)--(c3)--(c4)--(c1) (c1)--(c3) (c2)--(c4);
 \node[bag] (C1) at (6,1.5) {$\{1,2,3,4\}$};
 \node[sc] at (6,-22) {(c) 完全图：树宽$3$};
\end{scope}
\end{tikzpicture}
